% Transcription — fonds Alexandre Grothendieck, shelfmark 59, batch 2 (pages 21-27).
% Established from the facsimile archives/batches/59/batch-02.pdf (PDF page k =
% archivists' page 20+k, no cover sheet), cut from a copy of 59.pdf that was
% fetched through the project's own relay
% (https://grothendieck-relay.onrender.com/source/59.pdf), not uploaded by the
% user; tiles in archives/tiles/59/p21-p27. Per the skill transcribe-grothendieck.
% The folder is twenty-seven pages; this is the second and last batch (pages
% 1-20 are batch 1, transcribed separately and in parallel, so the notation
% here is fixed from these pages alone).
%
% Pass: Opus 5.5 (claude-opus-5-5), 2026-09-26 — first pass, unchecked against the pages by a human.
%
% Pages skipped: 21, a blank cream sheet carrying nothing but the faint
% show-through of the next leaf — no mathematics of his, skipped without
% description.
%
% Pages summarised: none. Pages transcribed: 22-27, six leaves in blue-black
% ink, his own pagination 1 to 6 (top right, pencil or ink). Page 24 (his 3)
% carries only two lines and the rest of the leaf is blank. Fast drafting:
% formulas carry, some connective words do not; interlinear insertions
% linked by strokes are placed where the stroke sends them, with a \note.
% His underlined words are set in \emph.
%
% What the batch is about. X/S proper, flat, of finite presentation, with
% f_*(O_X) = O_S universally and Pic^0_{X/S} an abelian scheme B; A =
% Pic^0_{B/S}. Pages 22-24: construction of the A-torsor Alb^1_{X/S} and of
% psi : X -> Alb^1_{X/S}, equivalently delta(x,y) = psi(y) - psi(x) with
% delta(x,x) = 0, given on sections by N_y N_x^{-1} for a Weil sheaf N on
% X x_S B; universal property for S-morphisms of X into torsors under abelian
% schemes (« ou plus généralement ... schémas en groupes commutatifs
% quelconques ? »); the inverse solution (Q^{-1}, psi^{-1}) and the resulting
% « convention de signe » for A = Pic^0_{B/S}, the first one being adopted.
% For X smooth of relative dimension 1: Pic^1_{X/S} is a solution of the
% Albanese problem via the diagonal, whence an isomorphism c : B -> A,
% c(psi(y) - psi(x)) = N_y N_x^{-1}; the open question whether c is the
% Theta-polarisation or (« sans doute ») its symmetric, reduced to a
% geometric question over an algebraically closed field, u(t) = class of
% t.Theta - Theta. Pages 25-27, « Schémas abéliens polarisés principaux »:
% for delta in NS_{A/S}, the torsor P^delta = Pic^delta_{A/S} under B and the
% sheaf L on A x P^delta; for delta a polarisation, pr_{2*}(L) locally free
% of rank chi(delta), the canonical divisor D_delta when chi(delta) = 1, its
% descent to Theta on P^delta; a Proposition in parts (a unlabelled), b), c),
% d): existence and uniqueness of (P, Theta) with deg Theta^g / g! = 1 for a
% polarisation of degree 1; Theta relatively ample; P = Pic^delta_{A/S}; for
% A = Alb^0 of a curve of genus g, P^delta = Pic^{g-1} = Alb^{g-1} with Theta
% going to W_{g-1}, and a boxed marginal « question de signe ! ». Then a
% framed and cancelled restart of the set-up (a symmetric divisorial
% correspondence C on X x_S X), and « (c bis) »: the isomorphism
% sigma_P : P -> P^{-1}, Theta^{-*} = sigma_P(Theta), a unique A-torsor
% isomorphism K between P^{-1} and P with K(Theta) = Theta^{-*}, K viewed as
% a section of P^{(x)_A 2}, and for the Jacobian case K comes from
% Omega^1_{X/S} under P^{(2)} = Pic^{2g-2}_{X/S}; a margin note asks to check
% that this is the canonical section for every delta in NS_{A/S}.
%
% Notation fixed once for the batch (from these pages; the Picard notation of
% the fonds, grepped from transcripts/): \underline{\mathrm{Pic}},
% \underline{\mathrm{Alb}}, \underline{\mathrm{NS}} — he underlines all three
% as representable functors; \mathcal{O}, \mathcal{L}, \Omega^1_{X/S}
% underlined as he writes it; \Theta for his barred or circled theta, which
% are the same letter throughout; \mathrm{cl} for « cl ». On page 25 the map
% A -> B defined by delta is written first with a tilde, later with a bar,
% and both are kept as written.
%
% Readings to check first: the interlinear words of p. 22 around « on prend »;
% « à faire » at the foot of p. 22; the layered corrections of the question
% on p. 23 (« Non, c'est sans doute la symétrique de la Theta-polarisation »);
% « correspondant » p. 23; the struck formula on p. 24; « donc d'un » and
% « a la vertu » on p. 25; the bracket « toute ... section de A qui invarie
% Theta est la section unité » on p. 26; the interlinear « comme A-torseur »
% of d); « ce qui le caractérise » at the head of p. 27; the letter of
% « (c bis) » and the superscript of Theta^{-*}; the cancelled frame and the
% margin note of p. 27.

\documentclass[11pt,a4paper]{article}
% Le préambule commun à toutes les transcriptions du fonds Grothendieck.
%
% Il tient en une idée : **l'incertitude se compose, elle ne se résout pas.**
% Une transcription qui lisse un mot douteux a détruit la seule chose pour
% laquelle on la faisait — un lecteur doit pouvoir savoir, à chaque endroit,
% ce qui était sur le feuillet et ce qui a été ajouté. Les six macros qui
% suivent sont donc l'appareil critique, et non de la décoration.
%
% Les mêmes commandes sont reconnues par `scripts/render.mjs`, qui produit la
% vue de lecture du site. Une macro ajoutée ici doit l'être aussi là-bas,
% faute de quoi le rendu échouera bruyamment — ce qui est voulu : un rendu
% silencieusement faux est pire qu'un rendu absent.

\ProvidesPackage{grothendieck}[2026/08/08 Transcriptions du fonds Grothendieck]

\usepackage{amsmath,amssymb,amsthm}
\usepackage{mathrsfs}
\usepackage{stmaryrd}
% Les diagrammes commutatifs sont partout dans ce fonds : c'est la langue
% même de la géométrie algébrique de Grothendieck, et une transcription qui
% les rendrait en prose ne transcrirait rien.
\usepackage{tikz-cd}
% Français par défaut — c'est la langue des feuillets — mais l'anglais est
% chargé aussi : la traduction et le résumé sont des documents anglais, et la
% typographie française (espaces avant les deux-points) y serait une faute.
% Ces documents basculent par \selectlanguage{english} en tête de corps.
\usepackage[english,french]{babel}
\usepackage{fontspec}
\usepackage{geometry}
\geometry{a4paper,margin=25mm,marginparwidth=28mm}
\usepackage{marginnote}
\usepackage{xcolor}
% ulem, not soul. `soul`'s \st and \ul are built on a character-by-character
% scan that XeLaTeX's font machinery breaks: the document fails with an
% undefined control sequence rather than degrading. ulem does the same two jobs
% and is Unicode-engine safe. `normalem` stops it hijacking \emph, which the
% transcriptions use for Grothendieck's own emphasis.
\usepackage[normalem]{ulem}
\usepackage{hyperref}
\hypersetup{colorlinks=true,linkcolor=black,urlcolor=black,pdfborder={0 0 0}}

% --- Métadonnées du lot -----------------------------------------------------
% Reprises telles quelles par le script de rendu, qui les affiche en tête de
% la vue de lecture. Elles fixent ce que le fichier prétend couvrir : sans
% elles, une transcription flotte sans point d'attache dans un fonds de seize
% mille feuillets.

\newcommand{\folder}[1]{\gdef\grothFolder{#1}}
\newcommand{\batch}[1]{\gdef\grothBatch{#1}}
\newcommand{\leaves}[2]{\gdef\grothFirst{#1}\gdef\grothLast{#2}}
\newcommand{\foldertitle}[1]{\gdef\grothFoldertitle{#1}}
\gdef\grothFolder{?}\gdef\grothBatch{?}\gdef\grothFirst{?}\gdef\grothLast{?}\gdef\grothFoldertitle{}

% --- Le résumé --------------------------------------------------------------
%
% Il ouvre la lecture modernisée, sous le titre « Résumé », et il est composé
% en petit. Ce n'est pas un ornement : c'est le seul passage du document qui
% ne soit pas des mathématiques, et un lecteur doit pouvoir le reconnaître
% avant de l'avoir lu — puis le sauter s'il connaît déjà le sujet, ou s'y
% arrêter s'il ne le connaît pas. Le filet qui le suit marque la reprise du
% corps.

\newenvironment{resume}
  {\small\setlength{\parskip}{0.45em}}
  {\par\medskip\hrule\medskip}

% --- Mots-clés ---------------------------------------------------------------
%
% En anglais, en clôture du résumé de la lecture modernisée : le vocabulaire
% moderne sous lequel chercher ce que les feuillets construisent. Le manifeste
% du site (`npm run manifest`) les extrait pour étiqueter la cote entière —
% cette ligne est la source unique des tags, il n'y a pas de fichier à côté.
\newcommand{\keywords}[1]{\par\smallskip\noindent
  {\footnotesize\textsc{Keywords} --- \textit{#1}}\par}

% --- L'appareil critique ----------------------------------------------------

% Le numéro de feuillet, en marge. C'est l'ancre qui permet de régler un
% désaccord sur une formule en regardant *un* feuillet plutôt qu'une chemise
% de six cents. Le site s'en sert aussi pour tourner les pages du fac-similé
% quand on fait défiler la transcription.
\newcommand{\leaf}[1]{%
  \marginnote{\footnotesize\color{gray!70}\textsf{#1}}%
  \label{leaf:#1}\ignorespaces}

% Illisible : on ne devine pas. Un mot inventé qui se lit comme les autres est
% le pire résultat possible.
\newcommand{\ill}{\textcolor{red!60!black}{[\,\ldots\,]}}

% Lecture douteuse : le mot est donné, et signalé comme incertain.
\newcommand{\uncertain}[1]{\uline{#1}}

% Ajout de l'éditeur — une lettre manquante, un mot sous-entendu.
\newcommand{\add}[1]{\textcolor{blue!50!black}{[#1]}}

% Biffé par Grothendieck lui-même. Ce qu'il a rayé fait partie du document :
% une rature dit souvent où la pensée a changé de direction.
\newcommand{\struck}[1]{\sout{#1}}

% Note du transcripteur, sur l'état matériel du feuillet.
\newcommand{\note}[1]{\par\smallskip{\small\color{gray!60!black}\textit{[#1]}}\par\smallskip}

% Note marginale de Grothendieck, distincte du corps du texte.
\newcommand{\marginal}[1]{\par\smallskip\noindent\hspace{1em}%
  {\small\color{gray!60!black}#1}\par\smallskip}

% --- Titre ------------------------------------------------------------------

\AtBeginDocument{%
  \begin{center}
    {\large\bfseries Fonds Alexandre Grothendieck --- cote n\textsuperscript{o}~\grothFolder}\\[2pt]
    {\itshape\grothFoldertitle}\\[4pt]
    {\small feuillets \grothFirst--\grothLast\ (lot \grothBatch)}\\[2pt]
    {\footnotesize\color{gray!60!black}Transcription établie d'après le fac-similé numérisé
     par l'Université de Montpellier.\\
     Les lectures incertaines sont soulignées, les illisibles notées [\,\ldots\,],
     les ajouts entre crochets.}
  \end{center}
  \vspace{1em}}

\endinput


\folder{59}
\batch{2}
\pages{21}{27}
\dating{[années 1960-1970]}
\watermark{Édition de démonstration}
\foldertitle{Construction des faisceaux inversibles sur schémas de Picard : notes manuscrites (s.d.)}

\begin{document}

\page{22}\note{paginé 1 par l'auteur}

Soit $X/S$ propre, plat, de présentation finie, $f_{*}(\mathcal{O}_X) \simeq
\mathcal{O}_S$ univ.\add{ersellement}, $\underline{\mathrm{Pic}}^{0}_{X/S}$ un
schéma abélien $B$. Soit $A = \underline{\mathrm{Pic}}^{0}_{B/S}$. On définit
alors un \struck{schéma} $A$-torseur $\underline{\mathrm{Alb}}^{1}_{X/S}$, et un
$S$-morphisme
\[
  \psi : X \longrightarrow \underline{\mathrm{Alb}}^{1}_{X/S} ,
\]
ou ce qui revient au même, un homomorphisme
\[
  \delta : X \times_S X \longrightarrow A \qquad \delta(x,y) = \psi(y) - \psi(x)
\]
tel que
\[
  \delta(x,x) = 0 .
\]
Pour ceci, on prend, \uncertain{étant données} deux sections $x, y$ de $X$ sur $S$,
\[
  \delta(x,y) = N_y N_x^{-1} \quad \text{dans } \underline{\mathrm{Pic}}^{0}_{B/S}(S)
\]
où $N$ est un module inversible de Weil sur $X \times_S B$ (qui existe grâce
à l'existence d'une section de $X$ sur $S$).\note{« étant données » est
écrit en interligne sous « on prend », et un trait le relie à la formule ; la
lecture en est douteuse}

On prouve la propriété universelle de $(Q = \underline{\mathrm{Alb}}^{1}_{X/S},
\psi)$ pour les $S$-morphismes de $X$ dans les \struck{schémas} torseurs sous
schémas abéliens (ou plus généralement, sous des schémas en groupes
commutatifs quelconques ?).

Si on définit $Q'$ par $\delta' = \delta^{-1}$, on trouve un isomorphisme
$Q' = Q^{-1}$, $\psi' = \psi^{-1}$, et $(Q', \psi')$ est une autre solution
au \struck{\ill{}} PB universel, canoniquement isomorphe au premier. Donc,
lorsqu'on définit $(\underline{\mathrm{Alb}}^{1}, \psi)$ comme solution du PB
universel, avec $A$ le schéma abélien des translations de
$\underline{\mathrm{Alb}}^{1}$, il y a une \emph{convention de signe}
\uncertain{à faire} pour définir un isomorphisme $A \simeq
\underline{\mathrm{Pic}}^{0}_{B/S}$. \emph{Nous prendrons la première
convention de signe}.

\page{23}\note{paginé 2 par l'auteur}

Supposons donc $X$ lisse de dim.\ relative $1$.\note{« lisse » est ajouté en
interligne ; un crochet en marge ouvre l'alinéa} On peut considérer alors
$\underline{\mathrm{Pic}}^{1}_{X/S}$, torseur sous $B =
\underline{\mathrm{Pic}}^{0}_{X/S}$, (correspondant aux faisceaux inv.\ sur
$X/S$ de degré relatif $1$.\note{« torseur sous $B = \ldots$ » est écrit en
interligne et renvoyé ici par un trait ; la parenthèse n'est pas fermée} Il
y a un morphisme canonique
\[
  X \xrightarrow{\ \psi\ } \underline{\mathrm{Pic}}^{1}_{X/S} ,
\]
\struck{et celui-ci} défini par la diagonale de $X \times_S X$, qui est un
diviseur relatif sur le deuxième facteur.

\begin{tikzcd}
X & X \times_S X \arrow[l, no head, "\mathrm{pr}_1"'] \arrow[d, no head, "\mathrm{pr}_2"] \\
S \arrow[u, no head] & X \arrow[l, no head]
\end{tikzcd}

\note{carré dessiné en marge gauche, à traits sans flèches}

Il fait de $\underline{\mathrm{Pic}}^{1}_{X/S}$ une solution du pb universel
d'Albanese. Il s'ensuit, grâce à la première convention de signe, un isom
\[
  c : B \xrightarrow{\ \sim\ } A = \underline{\mathrm{Pic}}^{0}_{B/S} ,
\]
qui s'explicite ainsi : si $y, x$ sont deux sections de $X$ sur $S$, alors
$\psi(y) - \psi(x)$ est une section de $\underline{\mathrm{Pic}}^{0}_{X/S}$
correspondant au diviseur relatif $y(S) - x(S)$, et on a
\[
  c(\psi(y) - \psi(x)) = c(\mathrm{cl}(y(S)) - \mathrm{cl}(x(S))) = N_y N_x^{-1}
\]
où $N$ est un faisceau de Weil sur $X \times_S B$.

\emph{Question à tirer au clair} : Cet isom $c$ définit-il la \struck{\ill{}}
$\Theta$-\emph{polarisation} sur $B$ ? \struck{(ou la symétrique \ill{})}
Non, c'est sans doute la \emph{symétrique} \struck{(\uncertain{polarisation}
\ill{})} de la $\Theta$-polarisation.\note{« Non, c'est sans doute la » est
écrit en interligne au-dessus de la parenthèse biffée ; « symétrique »,
souligné, surmonte une seconde parenthèse biffée : on donne l'état final de
la phrase, les biffures à leur place}

C'est une question « géométrique », dans laquelle on peut supposer $S$
spectre d'un corps alg.\ clos. On a une polarisation de $B$ par le diviseur
$\Theta$. L'homom \uncertain{correspondant} $u : B \to A$ est défini par
\[
  u(t) = \text{classe de } t.\Theta - \Theta .
\]
\note{« correspondant » est en interligne, relié par un trait ; un signe
empâté précède $t$ dans $u(t)$}
Si

\page{24}\note{paginé 3 par l'auteur ; le feuillet ne porte que ces deux
lignes}

$t$ est de la forme $\mathrm{cl}((y) - (x))$, on aura donc
\[
  u((y) - (x)) = \text{classe de } \ldots
\]
\struck{\ill{}}\note{à la suite, une formule en $\Theta$ est biffée d'un
trait épais, illisible sous la biffure}

\page{25}\note{paginé 4 par l'auteur}

\section*{Schémas abéliens polarisés principaux}

\note{titre souligné, de sa main, en tête de la page}

Soit $A/S$ un schéma abélien, muni d'une section $\delta$ de
$\underline{\mathrm{NS}}_{A/S}$, d'où un torseur $P^{\delta} =
\underline{\mathrm{Pic}}^{\delta}_{A/S}$ sous le schéma abélien dual $B =
\underline{\mathrm{Pic}}^{0}_{A/S}$, et un faisceau inversible $\mathcal{L}$
sur $A \times P^{\delta}$ \ldots\ D'autre part, $\delta$ définit
$\tilde{\delta} : A \to B$, et les opérations de $A$ sur $P^{\delta}$ via
$\tilde{\delta}$ sont celles déduites par transport de structure des
opérations de $A$ sur lui-même par translations.

Supposons que $\delta$ soit une polarisation, \struck{on} alors on sait que
les $H^i$ des $\mathcal{L}_x$ ($x \in P^{\delta}$) sont nuls pour $i > 0$, et
que \struck{\ill{}} $\mathrm{pr}_{2*}(\mathcal{L})$ est donc un module loc.\
libre de rang $\chi(\delta)$. Si $\chi(\delta) = 1$, c'est un module
inversible sur $P^{\delta}$, et il lui est associé un diviseur canonique
$D_{\delta}$ \struck{$\Theta$} sur $A \times P^{\delta}$ relativement à
$P^{\delta}$. D'autre part, $\bar{\delta} : A \to B$ est donc un isom., et
$P^{\delta}$ peut être regardé comme un $A$-torseur.

\struck{\ill{} d'un \ill{}} Le diviseur $D_{\delta}$ est stable par
opérations de $A$ opérant diagonalement sur $A \times P^{\delta}$.
\struck{Donc il} [Il provient \uncertain{donc d'un} diviseur sur le quotient
$(A \times P^{\delta})/A = P^{\delta}$, noté \struck{\ill{}} $\Theta$. Ce
dernier \uncertain{a la vertu} que $\chi(\Theta_{\delta}) = 1$.\note{le
crochet ouvert devant « Il provient » n'est pas refermé}

\emph{Proposition}. Soit $A$ un schéma abélien$/S$, de dim.\ relative $g$,
soit $\delta$ une polarisation sur $A$ de degré $1$. Alors il existe un
$A$-torseur\note{« de dim.\ relative $g$, soit » est ajouté en interligne et
renvoyé par un trait}

\page{26}\note{paginé 5 par l'auteur}

$P$, et un diviseur relatif positif $\Theta$ sur $P$, [tel que $\deg
\Theta^{g}/g! = 1$] et que la polarisation $\delta$ soit égale à celle
déduite de la polarisation de $P$ par le diviseur $\Theta$.\note{« relatif
positif » est en interligne ; la première partie de la proposition ne porte
pas de lettre a), les suivantes sont b), c), d)} \struck{\ill{}} Et
donnés $P'$ et $\Theta'$ ayant les mêmes vertus que $P$ et $\Theta$, il
existe un unique $A$-isom de $P$ sur $P'$, transformant $\Theta$ en
$\Theta'$ [en d'autres termes \uncertain{toute} \struck{\ill{}} section de
$A$ qui \uncertain{invarie} $\Theta$ est la section unité].

b) Si on a un torseur $P$ sous $A$, muni d'un diviseur relatif positif
$\Theta$ tel que $\deg \Theta^{g}/g! = 1$ \struck{de degré de polarisation
$1$}, alors $\Theta$ est ample rel$/S$ [et définit \struck{donc} une
polarisation de $P/S$ donc \struck{degré} de $A/S$. Cette dernière est de
degré de polarisation $1$.

c) Sous les conditions de b), considérons l'homom de composition $A \times P
\to P$, et soit $\Theta'$ l'image inverse de $\Theta$ dans $A \times P$.
Alors $\Theta$ est un diviseur positif relatif sur $A \times P/P$ et définit
donc un \struck{\ill{}} faisceau inversible sur $A \times P$, donc un homom
$P \to \underline{\mathrm{Pic}}_{A/S}$. Ce dernier est un isom de $P$ sur
$\underline{\mathrm{Pic}}^{\delta}_{A/S}$, où $\delta$ est la polarisation de
$A/S$ définie par \struck{\ill{}} $\Theta$ sur $P$.\note{« Alors $\Theta$ » :
sans prime sur la page, là où l'on attend $\Theta'$ ; « $/P$ » est ajouté en
interligne après $A \times P$}

d) Supposons que $A = \underline{\mathrm{Alb}}^{0}_{X/S}$, où $X/S$ est une
courbe propre et lisse relative de \emph{genre} $g$. \struck{Alors}
Considérons la polarisation canonique $\delta$ de $A$ \struck{définie par
$B$}, \struck{\ill{} \ill{}} le diviseur diagonal sur $X \times_S X$. Cette
polarisation est de degré $1$, et $P^{\delta}$ est canoniquement isomorphe,
\uncertain{comme} $A$-torseur, à $\underline{\mathrm{Pic}}^{g-1}_{X/S} =
\underline{\mathrm{Alb}}^{g-1}_{X/S}$, cet isomorphisme\note{« le diviseur
diagonal », « est de degré $1$ » et « comme $A$-torseur » sont écrits en
interligne, au-dessus de passages biffés ; l'ordre des mots est restitué
d'après les traits de renvoi}
\marginal{question de signe !}\note{encadré, en marge gauche au bas de la
page, relié par un trait à la dernière ligne}

\page{27}\note{paginé 6 par l'auteur}

transformant $\Theta$ en $W_{g-1}$ (\uncertain{ce qui le caractérise}).

\note{un trait horizontal suit. Vient ensuite un passage encadré, barré de
longs traits obliques, que l'on donne comme biffé}
\struck{Soit $X/S$ propre plat de prés.\ finie, $f_{*}(\mathcal{O}_X) \simeq
\mathcal{O}_S$ univ.\add{ersellement}, $\underline{\mathrm{Pic}}^{0}_{X/S}$ un
schéma abélien $B$, de dual $A = \underline{\mathrm{Alb}}^{0}_{X/S}$.
Supposons donnée sur \ill{} $(X \times_S X)$ une correspondance divisorielle
symétrique $C$, définissant donc une correspondance divisorielle symétrique
\ill{} $A$, donc une polarisation $\gamma$ de $A$, \uncertain{section} de
$\underline{\mathrm{NS}}_{A/S}$, soit $\gamma$.}\note{à l'intérieur du cadre,
les deux « symétrique » sont en interligne et plusieurs mots sont biffés
séparément avant la biffure d'ensemble}

\struck{\ill{} bis} (\uncertain{c} bis) Soit $(P, \Theta)$ comme dans b).
Considérons \struck{$P^{-1}$, $\Theta^{-1}$} $P^{-1}$, et l'isom (compatible
avec la \emph{symétrie} sur $A$)
\[
  P \xrightarrow{\ \sigma_P\ } P^{-1} ,
\]
et soit $\Theta^{-*} = \sigma_P(\Theta)$. Alors la polarisation définie par
$\Theta^{-*}$ sur $A$ est égale à celle définie par $\Theta$, donc il y a un
isomorphisme unique de $A$-torseurs
\[
  K : P \xleftarrow{\ \sim\ } P^{-1}
\]
tel que $K(\Theta) = \Theta^{-*}$. On peut aussi considérer $K$ comme une
section de $P^{\overset{A}{\otimes} 2}$
\[
  K \in \Gamma(S, P^{\overset{A}{\otimes} 2}) .
\]
\note{l'exposant de $\Theta^{-*}$ est lu douteusement : un signe moins suivi
d'une étoile ou d'un petit signe surchargé ; la lettre de « (c bis) » peut
se lire e. Dans $K : P \leftarrow P^{-1}$, le $P$ est surchargé}

Lorsque $A = \underline{\mathrm{Alb}}^{0}_{X/S}$, $X$, $\delta$ comme dans
d), donc $P \simeq \underline{\mathrm{Pic}}^{g-1}_{X/S}$ \struck{$\simeq
\underline{\mathrm{Alb}}$}, \struck{$K$ est la} moyennant l'isom.\ $P \simeq
\underline{\mathrm{Pic}}^{g-1}_{X/S}$, d'où $P^{(2)} \simeq
\underline{\mathrm{Pic}}^{2g-2}_{X/S}$, $K$ est la section provenant du
\struck{\ill{}} faisceau $\underline{\Omega}^{1}_{X/S}$ sur
$X$.\note{« $P \simeq$ » est ajouté en interligne}

\marginal{Vérifier que c'est la section \struck{\ill{}} canonique \ill{}
$(\underline{\mathrm{Pic}}^{\ill{}}_{A/S})$ \ill{} pour toute section $\delta$
de $\underline{\mathrm{NS}}_{A/S}$ \ldots}\note{en marge gauche, en oblique,
reliée par un long trait courbe à « Lorsque » et à la ligne de $K$}

\end{document}
